\documentclass[10pt,a4paper]{amsart}
\usepackage[margin=19mm]{geometry}
\usepackage{amsmath,amssymb,mathtools}
\usepackage[scr=rsfs,cal=euler]{mathalfa}
\usepackage{xfrac}
\usepackage[unicode,hidelinks,bookmarks=false]{hyperref}
\newcommand{\ie}{i.e.\ }
\newcommand{\cf}{cf.\ }
\newcommand{\resp}{respectively\ }
\newcommand{\Subsection}{\subsection}
\providecommand{\nobreakdash}{\nobreak\hbox{-}}
\setlength{\emergencystretch}{2em}
\multlinegap0pt
\usepackage{bm}
\usepackage{bbm}
\usepackage{dsfont}
\usepackage{xspace}
\usepackage{cancel}
\usepackage{mathtools}
\usepackage{amsthm}
\makeatletter
\@ifundefined{theorem}{\newtheorem{theorem}{Theorem}}{}
\@ifundefined{corollary}{\newtheorem{corollary}[theorem]{Corollary}}{}
\@ifundefined{lemma}{\newtheorem{lemma}[theorem]{Lemma}}{}
\makeatother
\newcommand {\Aut}{\mathrm{Aut}}
\newcommand {\Closure}{\mathrm{Closure}}
\newcommand {\Hom}{\mathcal{H}\kern -0.25ex{\mathit om}}
\newcommand {\ext}{\mathrm{ext}}
\newcommand{\Ho}{\operatorname{H}}
\newcommand{\Num}{\operatorname{Num}}
\newcommand{\h}{\operatorname{h}}
\title[Moduli spaces of stable bundles on ruled 3-folds and Brill-N]{Moduli spaces of stable bundles on ruled 3-folds and Brill-Noether problems}
\author{Laura Costa, Irene Mac\'ias Tarr\'io}
\date{Source: arXiv:2507.06882v1 (CC BY 4.0)}
\begin{document}
\maketitle

\noindent\emph{Source and licence.} Moduli spaces of stable bundles on ruled 3-folds and Brill-Noether problems. Version arXiv:2507.06882v1 (CC BY 4.0), distributed under the Creative Commons Attribution 4.0 International licence, as recorded in the arXiv licence metadata for this exact version and shown on the abstract page https://arxiv.org/abs/2507.06882v1. Full rights notice: licenses/arxiv-2507.06882-CC-BY-4.0.txt.\par\smallskip

\noindent\textit{Editorial scope.} Counted unit: the Theorem whose source label is \texttt{th\_irreducible\_comp} in the article (source0.tex lines 780--782, source label \texttt{th\_irreducible\_comp}), delivered together with the article's own complete original proof of it (source0.tex lines 783--857, one \texttt{proof} environment of 75 lines). No mathematics is written, repaired or re-derived: statement and proof are the article's own text line for line. Required context retained uncounted: cor\_inclusion (lines 511--514); lemma\_chib (lines 701--713). Every definition, display and label of those lines is retained unchanged. Omitted from this package and not redistributed: 1--510, 515--700, 714--779, 858--1191. Nothing inside a retained excerpt is omitted or altered: every retained body is delivered exactly as the article's lines are written, so no omission receipt of any kind accompanies this package. Citations: the retained excerpts carry no citation command at all, so no bibliography entry of the article is transcribed and no reference is redistributed. Reference adaptation: none was needed and none was made; every reference inside the delivered unit resolves inside it, so no unresolved reference or citation is produced. Printed slips kept literal: no printed word, symbol, hypothesis or number of the counted statement or of its proof was altered. Layout and class: the article's own class is \texttt{amsart} with packages including amsmath, tikz-cd, amssymb, amsthm, latexsym, amscd, xypic, ifthen, hyperref, graphicx, float, enumerate, xcolor, hyphenat; the wrapper is the lane's pinned \texttt{amsart} editorial wrapper at 10pt on A4, which restates the theorem-like environments the retained text needs and adds the article's own macro definitions that the retained text needs (\texttt{\textbackslash Aut}, \texttt{\textbackslash Closure}, \texttt{\textbackslash Ho}, \texttt{\textbackslash Hom}, \texttt{\textbackslash Num}, \texttt{\textbackslash ext}, \texttt{\textbackslash h}), with extra wrapper packages (bm, bbm, dsfont, xspace, cancel, mathtools). The delivered numbering is the unit's own. Assets: no retained span contains a figure environment, a graphics-inclusion command, a PDF-page inclusion, a table or a TikZ picture; no image file is copied into this package, the package declares no asset dependency and no image file is redistributed with it. Licence: version arXiv:2507.06882v1 of this article is distributed under the Creative Commons Attribution 4.0 International licence, as recorded in the arXiv licence metadata for this exact version and shown on the abstract page https://arxiv.org/abs/2507.06882v1. A full rights notice is delivered at licenses/arxiv-2507.06882-CC-BY-4.0.txt. Characters: every retained excerpt is pure ASCII, so the delivered engine, XeLaTeX with its default Latin Modern text font, reports no missing character.
\begin{corollary}
\label{cor_inclusion}
    Let $\xi\in\Num(X)\otimes\mathbb{R}$ defining a non empty wall $W^{\xi}$ and satisfying the $C-$condition. Assume that $W^{\xi}\cap \Closure(\mathcal{C})\neq\emptyset$ and $\xi\cdot L^{n-1}<0$ for  $L\in\mathcal{C}$.  Then, there exists an embedding $E_{\xi}(c_1,c_2)\hookrightarrow M_{\mathcal{C}}(2;c_1,c_2)$. %\textcolor{red}{revisar "some polarization in ", en lugar de "any"}
\end{corollary}

\begin{lemma}
\label{lemma_chib}
    Let $X=\mathbb{P}(\mathcal{E})\rightarrow\mathbb{P}^2$ be a ruled $3$-fold  over $\mathbb{P}^2$ and  consider $\xi_b= -S+2bH$ in $\Num(X)$ with $c_1'+1\leq b\in\mathbb{Z}$.
    Then:
    \begin{itemize}
        \item[i)] $\xi_b$ defines a non empty wall $W^{\xi_b}$ of type $(S, (b+1)SH-b^2H^2)$ and satisfies the $C-$condition.
        \item[ii)] The wall $W^{\xi_b}$ is only defined by $ \xi_b$ and $-\xi_b$, but $-\xi_b$ does not satisfy the $C-$condition. 
    \end{itemize}
    
    
    
    
\end{lemma}

 \begin{theorem}\label{th_irreducible_comp}Let $X=\mathbb{P}(\mathcal{E})\rightarrow\mathbb{P}^2$ be a ruled  $3$-fold  over $\mathbb{P}^2$ with $c'_i=c_i(\mathcal{E})$.  Let us fix $c_1'+1\leq b\in\mathbb{Z}$ and Chern classes $c_1=S$ and $c_2=(b+1)SH-b^2H^2$. % Let us consider $\xi_b=-S+2bH$. 
 Then,  $$M_{\mathcal{C}_b}(2;c_1,c_2)\neq\emptyset$$ and it contains an irreducible component of dimension $2b+\h^0\mathcal{E}-\h^0\mathcal{E}(-1)-1$.
        \end{theorem}     

        \begin{proof}
        By Lemma \ref{lemma_chib}, $\xi_b=-S+2bH$ defines the non empty wall $W^{\xi_b}$ of type $(c_1,c_2)$ and satisfies the $C-$condition. Thus, we can consider the family $E_{\xi_b}(c_1,c_2)$ of vector bundles given by a non trivial extension of type 
        \begin{equation}
        \label{echi}
         0\rightarrow\mathcal{O}_X(bH)\rightarrow E\rightarrow I_{[SH]}(S-bH)\rightarrow0.   
        \end{equation} Moreover, by Corollary \ref{cor_inclusion}, $E_{\xi_b}(c_1,c_2)$ is an irreducible component of the moduli space $M_{\mathcal{C}_b}(2;c_1,c_2)$. Since any $E\in E_{\xi}(c_1,c_2)$ is given by an extension of type (\ref{echi}), we get $$\dim{E_{\xi}}(c_1,c_2)=\ext^1 (I_Z(S-bH),\mathcal{O}_X(bH))+\dim\mathcal{L}-1,$$ where $\mathcal{L}$ parametrizes the codimension $2$ complete intersections of type $[Z]=SH$. 

Let us first compute $\ext^1 (I_{Z}(S-bH),\mathcal{O}_X(bH))$. Notice that $$\ext^1 (I_{Z}(S-bH),\mathcal{O}_X(bH))=\ext^2(\mathcal{O}_X(bH), I_Z(S-bH+K_X))=\h^2I_Z(S-2bH+K_X).$$

If we consider the exact sequence 
\begin{equation}
\label{eq_SH}
  0\rightarrow\mathcal{O}_X(-S-H)\rightarrow\mathcal{O}_X(-S)\oplus\mathcal{O}_X(-H)\rightarrow I_Z\rightarrow0  
\end{equation}

 twisted by $\mathcal{O}_X(S-2bH+K_X)$ and we take cohomology, we get 
\begin{equation}
\begin{array}{ll}
\label{suc_coh2}
\cdots &\rightarrow \Ho^2\mathcal{O}_X(-2bH+K_X) \oplus \Ho^2\mathcal{O}_X(S-(2b+1)H+K_X) \rightarrow \Ho^2I_Z(S-2bH+K_X) \rightarrow \\
&\rightarrow \Ho^{3}\mathcal{O}_X(-(2b+1)H+K_X) \rightarrow \Ho^{3}\mathcal{O}_X(-2bH+K_X) \oplus \Ho^{3}\mathcal{O}_X(S-(2b+1)H+K_X) \rightarrow \\
&\rightarrow \Ho^3I_Z(S-2bH+K_X) \rightarrow 0\\
\end{array}
\end{equation}



Notice that we have
\[
\begin{array}{llll}
    \h^2\mathcal{O}_X(-2bH+K_X) &= \h^1\mathcal{O}_X(2bH) &= \h^1\mathcal{O}_{\mathbb{P}^2}(2b) &= 0, \\
    \h^{3}\mathcal{O}_X(-(2b+1)H+K_X) &= \h^0\mathcal{O}_X((2b+1)H) &= \h^0\mathcal{O}_{\mathbb{P}^2}(2b+1) &= \binom{2b+2}{2}, \\
    \h^{3}\mathcal{O}_X(-2bH+K_X) &= \h^0\mathcal{O}_X(2bH) &= \h^0\mathcal{O}_{\mathbb{P}^2}(2b) &= \binom{2b+1}{2}, \\
    \h^2\mathcal{O}_X(S-(2b+1)H+K_X) &= 0, & & \\
    \h^{3}\mathcal{O}_X(S-(2b+1)H+K_X) &=0,&  & \\
    \h^{3}I_Z(S-2bH+K_X) &=\h^{3}\mathcal{O}_X(S-2bH+K_X)&= 0.  & \\
\end{array}
\]

Therefore, by (\ref{suc_coh2}) we obtain 

$\begin{array}{l}
\label{suc_coh}
0 \rightarrow \Ho^2I_Z(S-2bH+K_X) 
\rightarrow \Ho^{3}\mathcal{O}_X(-(2b+1)H+K_X) \rightarrow \Ho^{3}\mathcal{O}_X(-2bH+K_X) \rightarrow 0\\
\end{array}$

and thus 
\[
\begin{array}{rcl}
\h^2I_Z(S-2bH+K_X) &=& \h^{3}\mathcal{O}_X(-(2b+1)H+K_X) - \h^{3}\mathcal{O}_X(-2bH+K_X) \\
&=& \binom{2b+2}{2} - \binom{2b+1}{2} \\
&=& 2b+1.
\end{array}
\]


%Therefore, together with the fact that $\xi_b$ satisfies the $C-$condition and defines a non empty wall $W^{\xi_b}$,
Hence,   $E_{\xi_b}(c_1,c_2)\neq\emptyset$, which implies   $M_{\mathcal{C}}(2;c_1,c_2)\neq\emptyset$.
Finally, again using (\ref{eq_SH}), % by the exact sequence $$0\rightarrow\mathcal{O}_X(-S)\otimes\mathcal{O}_X(-H)\rightarrow\mathcal{O}_X(-S-H)\rightarrow I_{[SH]}\rightarrow0,$$
we get 
\[
\begin{array}{lll}
\dim\mathcal{L} &= \dim(\Hom(\mathcal{O}_X(-S-H),\mathcal{O}_X(-S)\oplus\mathcal{O}_X(-H)))-\dim\Aut(\mathcal{O}_X(-S-H))\\
&-\dim\Aut(\mathcal{O}_X(-S)\oplus \mathcal{O}_X(-H))+\dim I_f \\
&=  \h^0\mathcal{E}-\h^0\mathcal{E}(-1)-1,
\end{array}
\]
where $f\in\Hom(\mathcal{O}_X(-S-H),\mathcal{O}_X(-S)\oplus\mathcal{O}_X(-H))$ is a general element and $I_f$ denotes its isotropy group under the action of $\Aut(\mathcal{O}_X(-S-H))\times\Aut(\mathcal{O}_X(-S)\oplus\mathcal{O}_X(-H))$. Putting altogether, we obtain

$\begin{array}{rrl}
 \dim E_{\xi_b}&=&\ext^1 (I_Z(S-bH),\mathcal{O}_X(bH))+\dim\mathcal{L}-1    \\ &=&2b+1+\h^0\mathcal{E}-\h^0\mathcal{E}(-1)-1-1\\ 
     & =& 2b+\h^0\mathcal{E}-\h^0\mathcal{E}(-1)-1.
\end{array}$
 \end{proof}

\end{document}
